\documentclass[10pt]{article}

\usepackage[margin=1in]{geometry}
\usepackage[T1]{fontenc}
\usepackage{microtype}
\usepackage{booktabs}
\usepackage{amsmath}
\usepackage{xcolor}
\usepackage{enumitem}
\usepackage[hidelinks]{hyperref}

\definecolor{answered}{RGB}{27,120,55}
\definecolor{partial}{RGB}{184,108,0}
\definecolor{missing}{RGB}{176,35,35}
\newcommand{\answered}{\textcolor{answered}{\textbf{Answered from existing evidence}}}
\newcommand{\partialstatus}{\textcolor{partial}{\textbf{Partial---claim/correction ready}}}
\newcommand{\missing}{\textcolor{missing}{\textbf{Missing experiment/evidence}}}
\newcommand{\reviewer}[1]{\subsection{Reviewer #1}}
\newenvironment{reviewpoint}[2]
  {\subsubsection{#1}\noindent\textbf{Status:} #2\par\smallskip}
  {}

\title{ACCV 2026 Reviews: Rebuttal Answers and Lessons for WACV 2027}
\author{}
\date{29 August 2026}

\begin{document}
\maketitle

\section{Purpose and evidence policy}

This is a working research document, not yet a length-limited rebuttal. It records every
reviewer question, the answer currently supported by the submitted results, and the work
needed before reusing a claim in a WACV 2027 manuscript. Numerical claims below come from
the ACCV branch tables. Items marked \missing{} must not be answered with speculation or
presented as established conclusions.

The central lesson is that the ACCV submission over-referenced appendix material that was
not accessible with the main paper. Consequently, important evidence appeared to be
missing even when files existed in the repository. Every future claim must be supported by
a visible main-paper table/figure or by a correctly attached and resolvable appendix item.

\section{Short factual summary}

On the standard OccuFly test split, V-JEPA~2.1 FiLM-DPT-RGB obtains RMSE 4.621~m,
AbsRel 0.129, $\delta_1$ 0.833, and SILog 0.164. DepthAnything2-OccuFly obtains RMSE
4.844~m, AbsRel 0.134, $\delta_1$ 0.834, and SILog 0.112. Thus the submitted model is
better on RMSE, MAE, and AbsRel, essentially tied but slightly worse on the delta metrics,
and substantially worse on SILog. The defensible conclusion is metric-error
competitiveness, not overall superiority.

For SSC, the best submitted result is 37.68 SC IoU and 3.46 SSC mIoU. The best
Pose-MVS result is 32.55 SC IoU and 2.36 SSC mIoU, below the simpler RGB-context result.
These are negative but useful findings; they do not establish a new state of the art.

\section{Reviewer-by-reviewer answers}

\reviewer{LD5V}

\begin{reviewpoint}{LD5V-1: The paper lacks a clear end-to-end architecture/data-flow figure}{\partialstatus}
\textbf{Answer.} We agree. Figure~1 functions as an experiment map rather than an
end-to-end architecture diagram, and the remaining figures do not make the connection
among DTA, FiLM-DPT, Pose-MVS, depth-probability lifting, and SSC fusion sufficiently
concrete. The intended flow is:
\begin{align*}
&\text{RGB/video + calibration + altitude}
  \rightarrow \text{frozen V-JEPA 2.1}\\
&\quad\rightarrow \text{DPT/FiLM-DPT, DTA, or posed MVS branch}\\
&\quad\rightarrow \text{depth probabilities}
  \rightarrow \text{2D--3D lifting}\\
&\quad\rightarrow \text{SSC fusion and prediction head}.
\end{align*}
The revision should replace the screenshot-like overview with one legible diagram that
shows tensor/data flow, frozen versus trainable modules, and the mutually exclusive depth
branches. Each figure must be cited in the main text and checked at final PDF scale.
\end{reviewpoint}

\begin{reviewpoint}{LD5V-2: Table 3 caption/format is inconsistent and claimed baselines are absent}{\answered}
\textbf{Answer.} This is an authoring error. The caption claims MapAnythingV1.1,
Metric3Dv2, Depth Anything~2, and Depth Anything~3, but the body contains only the
DepthAnything2-OccuFly row. We will either restore all cited baseline rows using the same
OccuFly protocol or revise the caption to name only the included baseline. Best values
must be indicated per metric, not only for RMSE, and the discussion must cover all metrics.
\end{reviewpoint}

\begin{reviewpoint}{LD5V-3: Table 4 has one method and cannot establish interpolation}{\partialstatus}
\textbf{Answer.} We agree that a single row is not a comparison. The existing number only
shows that FiLM-DPT-RGB can be evaluated at held-out 40~m after training at 30/50~m; it
does not isolate FiLM, nor does it establish scene generalization. The submitted prose
must be weakened from ``supports physically meaningful conditioning'' to ``a preliminary
altitude-interpolation stress test.'' A valid replacement table requires an otherwise
identical unconditioned DPT control and preferably metadata ablations (altitude only,
intrinsics only, altitude+intrinsics).
\end{reviewpoint}

\begin{reviewpoint}{LD5V-4: Broken appendix references make evidence inaccessible}{\answered}
\textbf{Answer.} We agree and apologize. The repository contains altitude/scene MDE
tables and detailed SSC/per-class material, but the submitted package did not expose them
through valid cross-references. This is not merely cosmetic because the missing items were
the evidence for several claims. The revision must compile the main paper and appendix in
one deliverable (or attach the supplement correctly), resolve every reference, and reject
the build if ``??'' remains in the PDF or log.
\end{reviewpoint}

\begin{reviewpoint}{LD5V-5: FiLM is not isolated because plain DPT is unexpectedly weak}{\missing}
\textbf{Answer.} The current results do not isolate the FiLM contribution. A DPT result
that is worse than the linear probe is a warning of an optimization/configuration failure,
not a sound control from which to measure a large FiLM gain. Before making a causal FiLM
claim, we must audit preprocessing, frozen-layer selection, decoder initialization, depth
loss, learning-rate/weight-decay, output scaling/bins, and checkpoint selection; rerun DPT
and FiLM-DPT with identical settings and seeds; and add altitude-only and intrinsics-only
ablations. Until then the defensible statement is that the \emph{combined configured
FiLM-DPT system} performs best among the tested V-JEPA heads.
\end{reviewpoint}

\begin{reviewpoint}{LD5V-6: The altitude split is inconsistent and confounded by scene overlap}{\answered}
\textbf{Answer.} The experiment report identifies train at 30/50~m, validation at 40~m
scenes 01--04, and test at 40~m scenes 05--08. Therefore ``scenes 1--9'' is incorrect and
must be corrected to scenes 01--08 with the validation/test partition stated explicitly.
Because the same scene identities occur at other altitudes in training, the protocol holds
out altitude but not scene identity. We will not call it scene-generalization evidence. A
scene-disjoint altitude test is required for that claim.
\end{reviewpoint}

\begin{reviewpoint}{LD5V-7: Video does not account for camera-motion misalignment}{\missing}
\textbf{Answer.} Correct. DTA aggregates tubelet tokens without explicitly warping
features into the target camera frame. UAV ego-motion can therefore turn temporal
aggregation into spatial mixing. The negative/altitude-dependent video result must be
presented with this limitation. A controlled follow-up should compare no alignment,
pose-based feature warping, and (if available) optical-flow alignment using the same video
decoder and training protocol.
\end{reviewpoint}

\begin{reviewpoint}{LD5V-8: Pose-MVS is evaluated only through SSC, without direct depth results}{\missing}
\textbf{Answer.} We agree. Downstream SSC cannot diagnose whether Pose-MVS fails at
correspondence/depth estimation or during lifting/fusion. We must report direct Pose-MVS
depth metrics and qualitative depth/error maps on the same split as monocular depth,
including an oracle-depth lifting upper bound. Until then, Pose-MVS is an exploratory
negative SSC result, not evidence about MVS depth quality.
\end{reviewpoint}

\begin{reviewpoint}{LD5V-9: SSC mIoU is low and lacks per-class, oracle, and statistical analysis}{\partialstatus}
\textbf{Answer.} Per-class and altitude-wise outputs exist in the repository but were made
inaccessible by the appendix packaging failure. They must be restored and discussed,
especially class imbalance and the gap between SC IoU and SSC mIoU. Oracle-depth lifting
and multi-seed uncertainty are not present and must be added. The revised claim should be
diagnostic: better monocular depth does not automatically translate into stronger semantic
completion, and alignment between geometry, lifting, and context is a bottleneck.
\end{reviewpoint}

\reviewer{CWVJ}

\begin{reviewpoint}{CWVJ-1: There is no second frozen backbone}{\missing}
\textbf{Answer.} We agree that a single-backbone study cannot attribute observed behavior
specifically to V-JEPA~2.1. The current conclusion must be scoped to ``V-JEPA~2.1 under
this OccuFly protocol.'' For WACV, at least one alternative frozen encoder---DINOv3 is the
most directly motivated control raised by the reviewer---must use the same input
resolution, selected feature levels, decoder capacity, training budget, splits, and seeds.
If exact parity is impossible, parameter counts and deviations must be reported.
\end{reviewpoint}

\begin{reviewpoint}{CWVJ-2: Swapping V-JEPA into SSC barely changes performance}{\answered}
\textbf{Answer.} We accept this as a result rather than claiming a clear SSC gain. In the
submitted configurations, V-JEPA-derived priors move SSC modestly and inconsistently; the
architecture/lifting choice dominates. The contribution should therefore be framed as an
empirical boundary: strong frozen-feature MDE transfer does not by itself guarantee SSC
transfer.
\end{reviewpoint}

\begin{reviewpoint}{CWVJ-3: Plain DPT appears broken}{\missing}
\textbf{Answer.} This is the same validity issue as LD5V-5. We cannot use the present DPT
row as a clean causal baseline. The corrective experiment is a matched, multi-seed DPT
versus FiLM-DPT comparison after a training/configuration audit, with learning curves and
checkpoint-selection rules reported.
\end{reviewpoint}

\begin{reviewpoint}{CWVJ-4: Broken appendix references and unfinished manuscript}{\answered}
\textbf{Answer.} We agree. All evidence-bearing appendix items must ship with the paper,
all labels must resolve, and tables/figures must remain readable in the final PDF. A clean
compile check is part of the paper-release protocol below.
\end{reviewpoint}

\begin{reviewpoint}{CWVJ-5: Mixed comparison to DepthAnything2 and confounded interpolation}{\answered}
\textbf{Answer.} FiLM-DPT-RGB improves RMSE (4.621 vs. 4.844), MAE (3.523 vs. 4.193),
and AbsRel (0.129 vs. 0.134), while it is marginally worse on $\delta_1$ (0.833 vs.
0.834), $\delta_2$ (0.971 vs. 0.976), and $\delta_3$ (0.998 vs. 0.999), and markedly
worse on SILog (0.164 vs. 0.112). This pattern suggests competitive absolute metric
accuracy but less consistent scale-relative/log-space accuracy. It does not justify saying
the method ``beats DepthAnything2'' without qualification. The interpolation limitation is
addressed in LD5V-6.
\end{reviewpoint}

\begin{reviewpoint}{CWVJ-6: Single runs and potentially seed-sized margins}{\missing}
\textbf{Answer.} We agree. Key comparisons require at least three controlled seeds,
reported as mean $\pm$ standard deviation (and preferably paired per-scene differences or
bootstrap confidence intervals). Claims based on margins comparable to uncertainty must be
described as ties.
\end{reviewpoint}

\begin{reviewpoint}{CWVJ-7: UAV efficiency motivation conflicts with an eight-H100 evaluation}{\partialstatus}
\textbf{Answer.} The paper reports the DGX hardware but no training/inference cost. Camera
weight and power motivate visual sensing, but do not establish onboard deployability of a
V-JEPA-Gigantic pipeline. We must remove or narrow any deployment implication and report
trainable/total parameters, peak training memory, training GPU-hours, inference latency,
throughput, and memory at the evaluated resolution. If the model is not tested onboard,
state that explicitly.
\end{reviewpoint}

\begin{reviewpoint}{CWVJ-8: Missing related work}{\missing}
\textbf{Answer.} The related-work section must be expanded and the proposed experiments
positioned against: frozen multi-backbone 3D probing (Probe3D, Lexicon3D), calibrated or
intrinsics-aware depth (CAM-Convs and Metric3D's canonical-camera formulation), aerial
depth/geolocalization datasets (WildUAV, UseGeo, MidAir, AerialMetric), and real-flight
video/multi-elevation benchmarks (RuralScapes, Dronescapes, MESSI). Each citation and
dataset property must be verified from the primary paper/project before inclusion; no
unverified bibliographic entry should be added from memory.
\end{reviewpoint}

\begin{reviewpoint}{CWVJ-9: Contradictory SSC contribution wording}{\answered}
\textbf{Answer.} The method text says the goal is not a completely new SSC architecture,
while a contribution bullet can be read as claiming one. The consistent claim is that we
evaluate V-JEPA-derived geometry/context variants within CGFormer/FoundationSSC-style
pipelines, including an exploratory Pose-MVS branch. We will remove the new-architecture
claim unless novelty is clearly defined and validated against matching controls.
\end{reviewpoint}

\reviewer{y7qY}

\begin{reviewpoint}{y7qY-1: Explain why DepthAnything2 wins delta/SILog metrics}{\answered}
\textbf{Answer.} The exact comparison is given in CWVJ-5. RMSE/MAE emphasize absolute
metric errors (and RMSE strongly weights large errors), whereas delta accuracies measure
the fraction within multiplicative thresholds and SILog measures log-space consistency.
The observed trade-off is compatible with fewer/smaller large metric errors but a less
uniform multiplicative error distribution. This is an interpretation, not a demonstrated
mechanism; per-depth-range and error-distribution plots are needed to establish why.
\end{reviewpoint}

\begin{reviewpoint}{y7qY-2: SSC is weak and video/MVS is a negative contribution}{\answered}
\textbf{Answer.} We agree with the numerical characterization. The RGB-context
configuration is strongest; video/MVS does not improve it. We will present this as a
negative finding and investigate camera-motion alignment, direct MVS depth quality, and
oracle-depth lifting before making a positive video-geometry claim.
\end{reviewpoint}

\begin{reviewpoint}{y7qY-3: Broken references and figures not cited}{\answered}
\textbf{Answer.} We will restore the missing appendix in the distributed package, repair
all labels/references, cite every figure in the main text, and run the release checks listed
below. Claims whose support does not fit in the main paper will explicitly point to a
correctly attached supplementary item.
\end{reviewpoint}

\begin{reviewpoint}{y7qY-4: Figure text may violate minimum font size}{\partialstatus}
\textbf{Answer.} All figures must be regenerated as vector graphics where possible and
inspected at final single/double-column size. Labels will use at least the venue's minimum
font size. Screenshot-like tables should be replaced by native LaTeX tables or diagrams.
\end{reviewpoint}

\begin{reviewpoint}{y7qY-5: Equation 19 says scenes 1--8 while Table 4 says 1--9}{\answered}
\textbf{Answer.} Table~4 is erroneous. The recorded protocol uses 40~m scenes 01--04 for
validation and 05--08 for testing, hence scenes 01--08, not 01--09. We will correct the
table and state the overlap limitation.
\end{reviewpoint}

\section{Priority experiment plan before WACV 2027 claims}

\begin{enumerate}[leftmargin=*,label=\textbf{P\arabic*.}]
  \item \textbf{Repair the plain-DPT control and isolate conditioning:} matched DPT,
  altitude-only FiLM, intrinsics-only FiLM, and combined FiLM; same decoder, optimizer,
  budget, split, and at least three seeds.
  \item \textbf{Add a second frozen backbone:} DINOv3 (and optionally another strong
  encoder) under a capacity- and protocol-matched probe/decoder comparison.
  \item \textbf{Remove interpolation confounding:} create a scene-disjoint altitude
  protocol and compare FiLM with an unconditioned control.
  \item \textbf{Diagnose video geometry:} no alignment versus pose-warped/flow-aligned
  temporal features; report by altitude and motion magnitude.
  \item \textbf{Evaluate Pose-MVS directly:} standard depth metrics, qualitative errors,
  depth-bin calibration, and downstream SSC with predicted versus oracle depth.
  \item \textbf{Quantify uncertainty:} multi-seed mean/std and paired scene-level or
  bootstrap intervals for all headline comparisons.
  \item \textbf{Explain SSC failures:} accessible per-class and altitude tables, class
  frequencies, SC-versus-SSC analysis, and geometry/context/lifting controls.
  \item \textbf{Report computational cost:} parameters, GPU-hours, peak memory, latency,
  throughput, and a precise statement about whether onboard deployment was evaluated.
\end{enumerate}

\section{Paper-release checklist learned from ACCV}

\begin{itemize}[leftmargin=*]
  \item Build the exact upload artifact from a clean checkout; do not assume local
  appendix files are included.
  \item Search the compiled source/log/PDF for unresolved references, citations, ``??'',
  placeholder text, missing files, and overfull content.
  \item Cross-check every table caption against its actual rows, metrics, highlighting,
  split, units, and cited source.
  \item Ensure every figure/table is cited, legible at final size, and communicates a
  scientific point rather than serving as an experiment tracker.
  \item Keep a claim-to-evidence ledger: each abstract/contribution/conclusion claim maps
  to a main-paper or accessible supplementary table/figure.
  \item Distinguish observation, controlled ablation, and causal conclusion. Do not
  attribute an effect to FiLM, video, MVS, or V-JEPA without the corresponding control.
  \item Report negative results and uncertainty; do not promote a single-run numerical
  margin to a general conclusion.
\end{itemize}

\section{Candidate WACV 2027 framing}

The strongest defensible direction is not ``a new SSC architecture.'' It is a controlled,
multi-backbone study of when frozen visual features transfer from aerial metric depth to
semantic scene completion, and why that transfer can fail. The ACCV result supplies the
initial observation: frozen V-JEPA~2.1 plus a conditioned dense decoder is competitive in
absolute metric-depth errors, yet video/MVS and stronger depth do not automatically yield
better SSC. The new WACV experiments should turn that observation into evidence by adding
matched backbones, conditioning ablations, scene-disjoint generalization, motion-aware
video controls, direct MVS evaluation, oracle lifting, and uncertainty.

\end{document}
